\section{Nine suggested formulas of Kataoka}\label{sec:kataoka}
Kataoka's 2026 preprint \cite[\S6.3, Tables 6--7]{Kat} treats dyadic
towers with prescribed leading coefficient $\lambda$. Table 7 lists data
through level six and suggests a full formula for each $0\le\lambda\le8$.
The author verifies the coefficient of $n2^n$ there and leaves the
remaining coefficients unconfirmed. The following result establishes those
coefficients and the indicated ranges. Here $X,Y$ denote voltage
monomials for the two basis vectors, so $X^aY^b$ corresponds to the step
$(a,b)$. A superscript on $X$ is an exponent; a repeated entry means
parallel generators.

\Needspace{15\baselineskip}
\begin{theorem}\label{thm:kataoka}
Put $N=2^n$. For the nine step lists below, the valuation $T_n$ equals the
given expression for every $n$ in the displayed range.
\begin{center}\small
\begin{tabular}{cllc}
\toprule
$\lambda$ & Generators & $T_n$ & First $n$\\
\midrule
0 & $X^3Y^2,X,Y$ & $9N-2n-3$ & 3\\
1 & $X^2Y^2,X,Y$ & $nN+(15/2)N-4n-7$ & 3\\
2 & $X,Y,Y$ & $2nN+2n+1$ & 2\\
3 & $XY,X,Y$ & $3nN+9N-8n-9$ & 1\\
4 & $X^2,X,X,Y,Y$ & $4nN-4N+6n+8$ & 2\\
5 & $X^2Y^2,X^2Y,XY^2,X,Y$ & $5nN+5N-16n+22$ & 3\\
6 & $X^3,X,Y,Y$ & $6nN-10N+14n+14$ & 3\\
7 & $X^3Y^3,X^3Y^2,X^2Y^3,X^2Y,XY^2,XY$ & $7nN+4N-20n+14$ & 3\\
8 & $X^4,X,X,Y,Y$ & $8nN-16N+14n+6$ & 2\\
\bottomrule
\end{tabular}
\end{center}
\end{theorem}

\subsection{Bounded orders in the row $\lambda=0$}
For $u=(a,b)$, the projected steps are $3a+2b,a,b$. Reduce the expansion
of $F_u(1+t)$ modulo $(2,t^8)$, exactly as for the ternary table above.
For $(1,b)$, $0\le b\le7$, the orders are
$6,2,4,2,6,2,4,2$ in sequence. For $(a,1)$, $a=0,2,4,6$, they are all $2$.
These orders are stable on lifting each class modulo eight, and they are
strictly smaller than $e=2^{h-1}$ for $h\ge4$.
The sum of the twelve entries is $36$, so $W_h=36\cdot2^{h-3}$.
The first three exact shell sums $7,18,44$ give $T_3=63$. The recurrence
\eqref{eq:shells} yields $T_n=9\cdot2^n-2n-3$ for $n\ge3$.

\paragraph{Kataoka's row $\lambda=1$.}
Here the three steps are $(2,2),(1,0),(0,1)$. Put
$X=\zeta^a$, $Y=\zeta^b$, $Z=XY$, and write $\nu=e\,v_2$ at conductor
$2^h$, where $e=2^{h-1}$. The Laplacian eigenvalue has the exact factorization
\[
 F=4-(Z+1)H,\qquad
 H=(X+Y)/Z + (Z-1)^2(Z+1)/Z^2 .
\]
Indeed the third pair of translated terms adds $-(Z-Z^{-1})^2$ to the
square-torus eigenvalue. Work first with $h\ge4$.
For mixed parities the square-torus eigenvalue has order $2$ and
$\nu((Z-Z^{-1})^2)=4$. The $2e$ mixed orbits therefore contribute $4e$.

For odd $a,b$ take $a=1$. If $b\equiv-1\pmod4$, the first term of
$H$ has order $2$, since $\nu(X+Y)=\nu(1-\zeta^{b-1+e})=2$.
The second has order at least $12$ or vanishes. Thus $\nu(H)=2$.
Over these $e/2$ orbits, the exponents $1+b+e$ range over all multiples
of four modulo $2^h$. For a nonzero such exponent, say $k$, the factor
$Z+1$ has order $2^{v_2(k)}\le e$. The finite orders sum to
$\sum_{j=2}^{h-1}(e/2^j)2^j=(h-2)e$. One orbit has $Z=-1$, giving
$\nu(F)=2e$. On the others $\nu(F)=\nu(Z+1)+2<2e$. These orbits
together contribute $(h+1)e-2$.

For $b\equiv1\pmod4$ both $Z-1$ and $Z+1$ have order $2$.
Consequently the second term of $H$ has order $6$.
The first has order $4$ for $b\equiv5\pmod8$ and at least $8$
(including infinity) for $b\equiv1\pmod8$, since $8\mid e$.
The two contributions to $H$ have distinct orders in both cases.
The resulting orders of $F$ are $6$ and $8$, each occurring $e/4$
times. Their sum is $7e/2$. We have proved
\[
          W_h=(h+17/2)2^{h-1}-2\qquad(h\ge4).
\]
Direct cyclotomic reduction at $h=1,2,3$ gives the respective shell sums
$7,16,48$. Thus $T_1=5$, $T_2=19$, $T_3=65$. The recurrence
$T_h-T_{h-1}=W_h-2$ now gives
\[
             T_n=(n+15/2)2^n-4n-7\qquad(n\ge3).
\]

\subsection{The four even positive indices and the triangular lattice}
Let $A(U)=2-U-U^{-1}$. For $s\ge1$, put $e=2^{s-1}$ and take $\zeta$
of exact order $2^s$. The conductor shell consists of the pairs
$(X,Y)=(\zeta^a,\zeta^b)$ with at least one exponent odd modulo $2^s$.
If $W_s$ sums the valuations of the Laplacian eigenvalue $F$ over that
shell, the matrix-tree formula gives $T_s-T_{s-1}=W_s-2$.

For $\lambda=2m$, $1\le m\le4$, the eigenvalue is
$F=P_m(A(X))+2A(Y)$, where the integral polynomials satisfy
$P_m(0)=0$ and $P_m(z)\equiv z^m\pmod2$. Explicitly, the four
polynomials are $z$, $z(6-z)$, $z((3-z)^2+1)$, and
$z(4-z)[4-z(4-z)]+2z$. These identities follow from
$A(X^2)=A(X)(4-A(X))$ and $A(X^3)=A(X)(3-A(X))^2$.

Put $L=2^j$, the smallest power of two at least $m$.
Take $s\ge\max(3,j+2)$. If $a\ne0$ then
$t=v_2(A(X))=2\cdot 2^{v_2(a)}/e$. When $t\le1/L$, the odd
coefficient of $z^m$ gives $v_2(P_m(A(X)))=mt\le1$, since lower
coefficients are even and the constant term vanishes.
For odd $a$ this applies regardless of $b$; the $2e^2$ choices
contribute $4me$ to $W_s$. With $a$ even and $b$ odd, each stratum
$v_2(a)=k$ for $1\le k\le s-2-j$ contributes $2me$. There are
$e/2^k$ choices of $a$ and $e$ choices of $b$ in that stratum.
In the remaining $2L$ choices of $a$, either $a=0$, or $t\ge2/L$.
For the latter, the leading monomial has valuation
$mt\ge 2m/L>1+2/e$, and the terms of smaller positive degree
have valuation at least $1+t>1+2/e$. Higher terms also exceed this
bound. The strict bound for $2m/L$ follows at once for
$(m,L)=(1,1),(2,2),(3,4),(4,4)$ and the stated range of $s$.
The valuation of $F$ for all these $a$ is thus $1+2/e$.
In total,
\[
 W_s=\big[2m(s-j)+2L\big]e+4L.
\]

For $\lambda=3$ use the exact identity
$XYF=8XY-(1+X)(1+Y)(1+XY)$. For $s\ge2$, exactly one of $a,b,a+b$
is even. Label this residue $t$. The two odd arguments each give
$v_2(1+\zeta^{\mathrm{odd}})=1/e$. For each choice of the even position
and its value there are $e$ pairs. When $t=e$ the product is zero,
hence $v_2(F)=3$. For all the other even $t$ we have
$v_2(1+\zeta^t)\le1$, hence $v_2(F)=2/e+v_2(1+\zeta^t)<3$.
The sum of $v_2(1+\zeta^t)$ over those remaining even $t$ is $s-1$,
since $-\zeta^t$ takes all nonidentity values in $\mu_e$ once and
$\prod_{\xi\in\mu_e,\,\xi\ne1}(1-\xi)=e$. Thus
$W_s=3(s+4)e-6$.

For clarity, in the order $\lambda=2,4,6,8,3$, the shell totals at
$s=1,2,3$ are $(7,18,36)$, $(10,30,56)$, $(10,28,88)$,
$(10,28,80)$ and $(9,30,78)$, respectively. They can be obtained by
reducing each orbit's eigenvalue modulo $Z^e+1$ and taking the smallest
valuation in the basis $1,(Z-1),\ldots,(Z-1)^{e-1}$. The accompanying
integer checker performs these small computations independently.
Substituting in $T_0=0$, $T_s=T_{s-1}+W_s-2$, then using the shell
formulas, gives the corresponding five rows of Theorem~\ref{thm:kataoka}.

\subsection{Kataoka's two odd rows with a repeated diagonal factor}

The first list contains a basis, and the second contains $(2,1),(1,1)$, another basis. Thus all layers are connected.

Put $d=\lambda-2$, $m=2^s$, $e=2^{s-1}$ and let $F$ be the Laplacian symbol. Direct multiplication gives
\[
 \overline F\doteq (X+1)(Y+1)(XY+1)^d
 \quad\text{in }\mathbb F_2[X^{\pm1},Y^{\pm1}],
\]
where $\doteq$ means equality up to a Laurent monomial. For each primitive pair $(a,b)$ modulo $m$, exactly one of $a,b,a+b$ is even. Write $c$ for this even member and $W=\zeta^c$, $\zeta$ primitive of order $m$. Let $2^t$ be the order of $W$, including $t=0$ for $W=1$. For $t\ge1$ the order of $(1+z)^c-1$ in characteristic two is $B=e/2^{t-1}$, taking $c$ to be its residue in $[0,m)$. Replacing an exponent by a congruent one does not change the formal order below $m$. At the axis $a$ even or $b$ even the contact order is $B+d+1$. At the diagonal $a+b$ even it is $2+dB$. If the contact order is strictly below $e$, it equals the normalized valuation $e v_2(F(\zeta^a,\zeta^b))$ by the uniformizer test. For $s\ge5$ this settles both axes when $t\ge2$, and the diagonal when $t\ge L+1$, where $L=2$ for $d=3$ and $L=3$ for $d=5$.

For the remaining classes let $h=(\lambda-1)/2$ and let $Z$ denote a primitive root of order $m$. In an axis class expand $W^hZ^hF(W,Z)$ as $\sum_j c_j(Z-1)^j$; symmetry treats the second axis. In a diagonal class use $W^h ZF(Z,W/Z)$ instead. The following table gives the smallest pair $(v_2(c_j),j)$ in lexicographic order. A dash means that the generic calculation applies.
\[
\begin{array}{c|cc|cc}
&\multicolumn{2}{c|}{\text{axis}}&\multicolumn{2}{c}{\text{diagonal}}\\
 t&\lambda=5&\lambda=7&\lambda=5&\lambda=7\\\hline
 0&(1,2)&(1,4)&(2,2)&(2,2)\\
 1&(1,2)&(1,4)&(3,0)&(4,0)\\
 2&-&-&(2,0)&(3/2,2)\\
 3&-&-&-&(2,0)
\end{array}
\]
To verify the entries, at $W=1,-1$ respectively the axis coefficient lists, in increasing $j$, are
\[
\begin{array}{c|ll}
5&(0,0,-10,-10,-2)&(8,16,6,-2,0)\\
7&(0,0,-28,-56,-42,-14,-2)&(-16,-48,-56,-32,-10,-2,0).
\end{array}
\]
For the diagonal the lists in increasing $t$ are
\[
\begin{array}{c|l}
5&(0,0,-4),\quad(8,8,0),\quad(-12,-12,0)\\
7&(0,0,-4),\quad(-16,-16,0),\quad
 (-16W,-16W-4,-2W-2),\quad(12W^3,12W^3,0).
\end{array}
\]
These expansions use only the identity $\Phi_{2^t}(W)=0$ and $\Phi_1(W)=W-1$ when $t=0$. For $W$ of order four we have $v_2(1+W)=1/2$, explaining the fractional entry. If the table entry is $(H,J)$, the normalized valuation is $eH+J$: the coefficient valuations lie in $\varphi(2^t)^{-1}\mathbb Z$. For $s\ge5$, the polynomial degree (at most six on the axes and two on the diagonal) is smaller than $e/\varphi(2^t)$ in every tabulated case. Thus the minimum term is unique.

To sum each class over its primitive pairs, for each fixed $W$ there are $e$ choices for the primitive coordinate $Z$. A normalized order $D$ therefore contributes $D$ for each such $W$, and there are $\varphi(2^t)$ choices of exact order $2^t$. In a single axis class the generic sum is $(s-2)e+(d+1)(e-2)$, while its two exceptional classes contribute $2(e+d-1)$. Each axis total is consequently $(s+d+1)e-4$. The generic diagonal sum is
\[
 d(s-L-1)e+2(e-2^L).
\]
Its exceptional classes contribute $9e+2$ for $d=3$ and $17e+6$ for $d=5$. Adding all three classes yields
\[
 W_s=(5s+10)e-14\quad (\lambda=5),\qquad
 W_s=(7s+11)e-18\quad (\lambda=7),\qquad s\ge5.
\]
The four initial valuations of the trees are $(8,34,134,358)$ and $(10,50,154,446)$ respectively. They can be verified by the finite cyclotomic basis calculation given above; the accompanying script recomputes all twelve initial layers. Since $T_s-T_{s-1}=W_s-2$, summation proves the asserted formulas.

